% Sustituye el contenido de ejemplo por el registro de tu propia sesión.
\section{\lblGlance}
\begin{table}[H]
  \centering
  \begin{tabularx}{\linewidth}{@{}l@{\hspace{1.2cm}}X@{}}
    \toprule
    \metalabel{\lblDate} & \RepDate \\
    \metalabel{\lblDuration} & \RepDuration \\
    \metalabel{\lblAttendees} & Dos miembros del equipo de soporte y la responsable de operaciones \\
    \metalabel{\lblFormat} & Videollamada con pantalla compartida \\
    \metalabel{\lblTopics} & Agentes de IA en atención al cliente, revisión de respuestas, manejo de datos de clientes \\
    \bottomrule
  \end{tabularx}
\end{table}

\begin{itemize}
  \item Los agentes ya pueden redactar respuestas para el equipo de soporte, pero no deben enviarlas solos.
  \item Los datos de clientes necesitan un límite claro antes de que un agente lea cualquier ticket.
  \item Un piloto pequeño con resultados medidos es mejor punto de partida que un despliegue completo.
\end{itemize}

\section{\lblQuestions}
\begin{enumerate}
  \item ¿Qué tareas de soporte puede hacer hoy un agente de IA de forma segura?
  \item ¿Cómo evitamos que un agente comparta datos de clientes que no debería ver?
  \item ¿Quién revisa el trabajo del agente y cuánto tiempo lleva?
  \item ¿Cómo sabremos que el agente ayuda y no solo que es rápido?
\end{enumerate}

\section{\lblFound}
\subsection{Los tickets mezclan solicitudes rutinarias y sensibles}
\findingpart{\lblContext}
El equipo atiende unos 300 tickets por semana. La mayoría son consultas sobre
pedidos y contraseñas; unos pocos tratan reembolsos, quejas o datos personales.

\findingpart{\lblObservation}
Los tickets rutinarios encajan bien con un agente. Los sensibles están en la
misma cola sin etiqueta, así que un agente no puede distinguirlos.

\subsection{Las respuestas no se revisan de forma consistente}
\findingpart{\lblContext}
Hoy cada persona del equipo revisa las respuestas a su manera, y no hay un
estándar escrito de tono ni de exactitud.

\findingpart{\lblObservation}
Sin un estándar compartido, quien revisa no puede decir si un borrador es
bueno, y el equipo no puede medir al agente contra nada.

\begin{insight}
  El límite seguro de un agente lo fija lo que puede ver y lo que puede enviar,
  no lo inteligente que sea. Conviene fijar primero esos dos límites.
\end{insight}

\begin{mermaid}[style=editorial,width=0.9\linewidth,
    caption={Manejo propuesto de un ticket},pos=H]
  flowchart LR
  A[Ticket] --> B[El agente redacta] --> C[Una persona revisa] --> D[Envío]
\end{mermaid}

\section{\lblRecommendations}
\begin{table}[H]
  \centering
  \begin{tabularx}{\linewidth}{@{}lXll@{}}
    \toprule
    & \sffamily\small\bfseries\lblRecommendation & \sffamily\small\bfseries\lblPriority &
    \sffamily\small\bfseries\lblEffort \\
    \midrule
    \recrow{1}{Mantener a una persona entre el agente y el cliente}{Solo borradores; una persona envía cada respuesta.}{high}{small}
    \recrow{2}{Ocultar los datos personales al agente}{Enmascarar nombres, correos y datos de pago antes de que lea el ticket.}{high}{medium}
    \recrow{3}{Escribir un estándar de respuestas de una página}{Da a quienes revisan y al piloto una medida común de calidad.}{medium}{small}
    \recrow{4}{Hacer un piloto de cuatro semanas con tickets rutinarios}{Medir el tiempo ahorrado y los errores antes de ampliar el alcance.}{medium}{medium}
    \recrow{5}{Añadir una revisión mensual de los errores del agente}{Mantiene las reglas al día cuando cambian los tickets.}{low}{small}
    \bottomrule
  \end{tabularx}
\end{table}

\section{\lblResources}
\begin{table}[H]
  \centering
  \begin{tabularx}{\linewidth}{@{}Xl@{}}
    \toprule
    \sffamily\small\bfseries\lblResource & \sffamily\small\bfseries\lblLink \\
    \midrule
    \resrow{Lista de control para revisar las respuestas de un agente}{https://example.org/checklist}
    \resrow{Plantilla de estándar de respuestas de una página}{https://example.org/reply-standard}
    \resrow{Guía para enmascarar datos personales en tickets}{https://example.org/masking}
    \resrow{Hoja de trabajo del plan piloto}{https://example.org/pilot-plan}
    \bottomrule
  \end{tabularx}
\end{table}

\section{\lblNextSteps}
\begin{table}[H]
  \centering
  \begin{tabularx}{\linewidth}{@{}Xlr@{}}
    \toprule
    \sffamily\small\bfseries\lblAction & \sffamily\small\bfseries\lblOwner &
    \sffamily\small\bfseries\lblTarget \\
    \midrule
    \actrow{Etiquetar los tickets sensibles en la cola}{client}{20 de noviembre de 2026}
    \actrow{Enviar la guía de enmascarado y un ejemplo de configuración}{consultant}{23 de noviembre de 2026}
    \actrow{Redactar el estándar de respuestas de una página}{client}{27 de noviembre de 2026}
    \actrow{Revisar el borrador del estándar y el plan piloto}{consultant}{2 de diciembre de 2026}
    \actrow{Iniciar el piloto de cuatro semanas}{client}{7 de diciembre de 2026}
    \bottomrule
  \end{tabularx}
\end{table}

\section{\lblNotes}
\begin{rawnotes}
  \begin{itemize}
    \item Al equipo le preocupan las promesas de reembolso erróneas; el agente nunca debe afirmar un reembolso.
    \item La herramienta actual guarda el historial completo, incluidos dígitos de tarjeta en texto libre.
    \item La responsable de operaciones quiere un número semanal: tickets respondidos, ediciones por borrador.
    \item Idea: empezar solo con estado de pedidos y añadir ayuda con contraseñas después.
  \end{itemize}
  Esquema de reglas acordado en la llamada:
  \begin{code}[lang=yaml]
agent:
  can_read: [order_status, shipping]
  can_send: false
  blocked: [refund, complaint, payment]
  \end{code}
\end{rawnotes}

\section{\lblFollowUp}
{\raggedright\hyphenpenalty=10000\exhyphenpenalty=10000 \RepBooking\par}
